%=====================================================================
\chapter{Sixteen-Week Schedule and Assessment Map}\label{app:schedule}
%=====================================================================

A mapping from the 33 chapters onto a standard semester of three contact hours
per week, with the assessment points that \reg{PhD Regs 2024}{\S51} and
\reg{MPhil/MS Regs 2024}{\S36} require. Adjust it to your own calendar; the
dependencies between chapters are what matter, not the week numbers.

\begin{regbox}[Semester shape --- PhD \S53, \S41; MPhil/MS \S15, \S17]
The semester is 18 weeks, of which \textbf{2 are reserved for the mid-term and
final examinations}, leaving 14 to 16 weeks of teaching. This schedule assumes
16. A contingency for 14 is given at the end of this appendix.

\medskip
This is a 3(3--0) course: three credit hours of theory, three contact hours of
class work per week, no laboratory hour. Practical work at this level is
expected to be accomplished through the research itself rather than through
scheduled labs (\S51.2(v)).
\end{regbox}

\section{The Teaching Schedule}

\begin{center}
\small
\begin{longtable}{@{}L{1.2cm}L{4.0cm}L{5.0cm}L{3.9cm}@{}}
\toprule
\textbf{Wk} & \textbf{Chapters} & \textbf{In class} & \textbf{Portfolio due} \\
\midrule
\endfirsthead
\toprule
\textbf{Wk} & \textbf{Chapters} & \textbf{In class} & \textbf{Portfolio due} \\
\midrule
\endhead
\bottomrule
\endfoot
1 & 1 Contribution\newline 2 Philosophy
  & What a contribution is; the paradigms behind method choice
  & --- \\
2 & 3 Questions\newline 4 Publication system
  & FINER and PICOC on the students' own topics; HJRS categories
  & \textbf{P1} Topic statement \\
3 & 5 SLR
  & Protocol workshop: search strings built and run live
  & --- \\
4 & 6 Grey literature\newline 7 Bibliometrics
  & Snowballing; Zotero and VOSviewer at the keyboard
  & \textbf{P2} SLR protocol \newline \emph{Quiz 1} \\
5 & 8 Appraisal\newline 9 Design science
  & Three-pass reading of one assigned paper; DSR artefact types
  & --- \\
6 & 10 Experiments\newline 11 Quasi-experiments
  & Designing a code-review experiment; confounding in repository data
  & --- \\
7 & 12 Case study\newline 13 Survey
  & Unit of analysis; instrument critique
  & \textbf{P3} Design choice memo \\
8 & \textbf{Mid-term}\newline 14 Qualitative
  & \textbf{Mid-term examination, Chapters 1--13}; then coding practice
  & --- \\
9 & 15 Mixed methods\newline 16 Benchmarking
  & Joint displays; workload characterisation
  & --- \\
10 & 17 Measurement\newline 18 Sampling and power
  & Operationalising a construct; a power calculation each
  & \textbf{P4} Draft instrument \\
11 & 19 Inference
  & The p-value misinterpretation clinic; effect sizes in R
  & \emph{Quiz 2} \\
12 & 20 Regression\newline 21 SEM
  & Model interpretation; a PLS-SEM run in \texttt{seminr}
  & --- \\
13 & 22 ML evaluation\newline 23 Qualitative analysis
  & A leakage audit on a supplied notebook; coding reliability
  & \textbf{P5} Analysis plan \\
14 & 24 Validity\newline 25 Reproducibility
  & Writing a threats section; packaging an artefact
  & \emph{Quiz 3} \\
15 & 26 Pre-registration\newline 27 Ethics\newline 28 Integrity\newline 29 Generative AI
  & The rigour and integrity block; similarity report interpretation
  & \textbf{P6} Pre-registration \\
16 & 30 Synopsis\newline 31 Writing\newline 32 Publishing\newline 33 Defending
  & Production workshop and defence-style presentations
  & \textbf{P7} Synopsis draft\newline \textbf{Presentation} \\
\end{longtable}
\end{center}

\section{The Portfolio and How It Maps onto the Regulated Weights}

The seven portfolio pieces are one continuous piece of work: each is a section
of the synopsis you will actually submit. They run \emph{inside} the weights
fixed by regulation, not on top of them.

\begin{center}
\small
\begin{tabular}{@{}L{4.8cm}cL{8.2cm}@{}}
\toprule
\textbf{Regulated component} & \textbf{Weight} & \textbf{How it is earned in this course} \\
\midrule
Participation / group work & 5\%  & Weekly seminar; the week 5 three-pass
                                    reading and the week 11 clinic are the
                                    two graded discussions \\
Quiz                       & 5\%  & Three surprise quizzes (weeks 4, 11, 14),
                                    5 questions each, as \S51.1(b) prescribes \\
Assignment                 & 5\%  & The portfolio, P1--P7, marked as one
                                    submission at the end \\
Presentation / seminar     & 5\%  & Week 16, 15 minutes, run as a mock DRC
                                    synopsis defence with questions \\
Mid-term examination       & 30\% & Week 8, Chapters 1--13, 30 marks in
                                    1~h~15 \\
Final-term examination     & 50\% & Whole syllabus, 50 marks; 20--30\% of the
                                    weight on pre-midterm material \\
\bottomrule
\end{tabular}
\end{center}

\begin{alertbox}[The portfolio is worth more than 5\%, whatever the table says]
Read that table honestly and it looks perverse: the largest piece of work in
the course carries the smallest weight, because the regulation caps assignments
at 5\% and cannot be varied by the instructor.

\medskip
It is not perverse, because of where the essay questions come from. The
mid-term essay and the two-to-four final essay questions are drawn from the
same material the portfolio makes you work through --- justify a design
against alternatives, state the threats to validity of a given study, write a
search string for a stated question. A student who has done P1--P7 properly
has already rehearsed 80\% of the marks. A student who has not will be writing
those essays from memory of the lectures, which is a much worse position.
\end{alertbox}

\subsection*{The seven portfolio pieces}

\begin{center}
\small
\begin{tabular}{@{}cL{4.0cm}L{9.3cm}@{}}
\toprule
 & \textbf{Piece} & \textbf{What it must contain} \\
\midrule
P1 & Topic statement & The problem, the gap, and why it matters here.
     Template C1 and C2. One page. \\
P2 & SLR protocol & Research questions, PICOC, search strings per database,
     inclusion and exclusion criteria, quality checklist, extraction form.
     Templates C3 and C4. \\
P3 & Design choice memo & The design you have chosen, two you rejected, and
     the constraint that decided it. Two pages. \\
P4 & Draft instrument & The questionnaire, interview guide, experiment
     protocol or extraction sheet your design needs, with each item traced to
     a construct. Templates C7--C10. \\
P5 & Analysis plan & The test or model per research question, the effect size
     you will report, the assumption checks, and the multiple-comparison
     correction. Written \emph{before} any data exists. \\
P6 & Pre-registration & P3--P5 assembled into the template of \chref{prereg},
     with the deviations policy stated. Template C12. \\
P7 & Synopsis draft & The IUB synopsis structure, populated. Template C14.
     This is the deliverable that outlives the course. \\
\bottomrule
\end{tabular}
\end{center}

\section{Two Tracks, One Course}

The teaching is common; the portfolio target differs, because the two degrees
demand different things of you (Appendix~\ref{app:compliance}).

\begin{center}
\small
\begin{tabular}{@{}L{3.4cm}L{5.8cm}L{5.8cm}@{}}
\toprule
 & \textbf{MS IT} & \textbf{PhD IT} \\
\midrule
P7 targets
  & A thesis synopsis deliverable within two semesters of research time
  & A dissertation synopsis with a first publishable increment identified \\
P2 scope
  & A focused review, roughly 25--40 primary studies
  & A review scoped to stand alone as a submittable paper \\
P6 emphasis
  & Feasibility: can this be finished and defended in time?
  & Publication plan: which venue, which category, and when relative to
    synopsis approval (\S15) \\
Extra reading
  & \chref{casestudy}, \chref{survey}
  & \chref{slr}, \chref{designscience}, \chref{publishing} \\
\bottomrule
\end{tabular}
\end{center}

\begin{conceptbox}[A note for PhD scholars on the comprehensive examination]
This course is examined in its own right, but its content reappears. The PhD
comprehensive examination (\S11) is a written and oral test of doctoral-level
mastery, set by the teachers who taught the courses and marked with an external
examiner, at 65\% to pass. Methodology questions in that examination are drawn
from exactly this syllabus. Keep the portfolio; it is the best revision
material you will have.
\end{conceptbox}

\section{If You Have Fourteen Teaching Weeks, Not Sixteen}

\reg{PhD Regs 2024}{\S53} permits as few as 14. Compress as follows, and do not
compress anything else:

\begin{tight}
  \item Merge weeks 9 and 10: \chref{mixed} and \chref{benchmarking} become
    assigned reading with a checkpoint, and \chref{measurement} and
    \chref{sampling} take the class time. Both merged chapters remain
    examinable.
  \item Merge weeks 12 and 13: teach \chref{regression} and \chref{mleval} in
    class; \chref{sem} becomes a guided exercise for those whose design needs
    it, and \chref{qda} likewise for those doing qualitative work. Most
    students need one or the other, rarely both.
\end{tight}

\noindent
Do not compress \chref{slr} (week 3), the integrity block (week 15) or the
production workshop (week 16). The first is where the portfolio is actually
built, and the last two are where students learn the rules they can otherwise
fail on without ever having been taught them.

\section{For the Course File}

\reg{PhD Regs 2024}{\S49} makes a course file compulsory and lists its
contents. This appendix supplies several of them directly:

\begin{center}
\small
\begin{tabular}{@{}L{7.0cm}L{7.3cm}@{}}
\toprule
\textbf{Required item} & \textbf{Where it comes from} \\
\midrule
Description of the course        & The front matter of this volume \\
Course coding                    & From the departmental scheme of study \\
Weekly teaching schedule         & The table in \S1 of this appendix \\
Copy of material / outline distributed & This volume \\
Topic and evaluation record of assignments & The portfolio table, \S2 \\
Copy of the quiz / test          & Instructor's own; three per \S2 \\
Mid and final papers with key    & Instructor's own; pattern in
                                   Appendix~\ref{app:compliance} \\
Course award list                & Instructor's own \\
\bottomrule
\end{tabular}
\end{center}

\begin{conceptbox}[Two assessment design notes for the instructor]
\textbf{Set at least one essay question that cannot be answered by
recall.} ``Describe the PRISMA flow diagram'' tests memory. ``A student
proposes to answer RQ1 with a survey of 40 developers recruited from one
company. State the strongest threat to validity and the cheapest design change
that would reduce it'' tests the thing this course is for.

\medskip
\textbf{Give credit in the portfolio for a design that is honestly abandoned.}
P3 asks for two rejected designs and the constraint that decided it. A student
who discovers in week 10 that their power calculation makes the planned
experiment infeasible, and says so, has learned more than one who proceeds
regardless. Mark that as a success, not a failure --- it is precisely the
judgement \S14.3(C) asks examiners to look for in a dissertation.
\end{conceptbox}
